\chapter{Implementazione}

\section{Sviluppo dell'Importer Infor Exchange}

\subsection{Gestione dei Formati e Parsing}
L'importer è stato sviluppato come un modulo d'input che si interpone tra i file d'ingresso
(prodotti dai vecchi server o estratti tramite tool legacy, a volte anche Access) e le tabelle di definizione del nuovo sistema.

La libreria dedicata alla gestione dei metadata (\texttt{txamgdll0011}) prende in carico l'interpretazione dei metadati
(quali formati delimitati o logiche ASCII a lunghezza fissa) eseguendo conversioni su base parametrica secondo
l'enumerato \texttt{txamg.fast}.

Ogni riga recuperata viene controllata per garantire la conformità delle codifiche testo e i troncamenti (\textit{trim}) o padding dei campi vuoti.

\begin{listing}[H]
\begin{minted}[frame=lines, framesep=2mm, baselinestretch=1.2, fontsize=\footnotesize, linenos]{c}
function void GeneraImportUnpack(domain txamg.dfil datafile)
{
    |* Generazione dinamica della scansione della riga ASCII
    |* print_format conterrà una stringa tipo "|%s|%s|%s|"
    |* string_list conterrà la lista delle variabili di destinazione

    select txamg111.* from txamg111
    where txamg111._index1 = {:txamg110.grou, :txamg110.dfil}
    order by txamg111.fdno
    selectdo
        print_format = print_format & first.pipe & "%s|"
        string_list = string_list & comma & "imp.str."
                      & trim$(txamg111.fdnm)
        comma = ","
        first.pipe = ""
    endselect

    Riga("function extern long unpack."&datafile&"(...) {")
    Riga("    ret = string.scan(buffer_"
              &trim$(datafile)&","""
              &trim$(print_format)&""","&trim$(string_list)&")")
    Riga("    ...")
    Riga("}")
}
\end{minted}
\caption{Esempio di generazione del codice per il parsing di stringhe ASCII (da \texttt{txamgdll0012}).}
\label{lst:parsing_gen}
\end{listing}

In questo estratto viene mostrato come il generatore iteri sui campi definiti per
il \textit{datafile} (\texttt{txamg111}) per costruire una chiamata a \texttt{string.scan}.
Questa funzione Baan C viene poi usata nello script generato per "spacchettare" la riga letta dal file ASCII nelle
variabili corrispondenti ai campi.

\subsection{Ricostruzione delle Definizioni dei File da Migrare}
L'obiettivo secondario dell'importer è auto-compilare il setup dello stesso \textit{Migration Tool}.
Utilizzando le definizioni originarie dei file (le tabelle legacy), il programma crea dinamicamente o aggiorna una
mappatura interna dei \textit{Data Files} attraverso Sessioni back-end come \texttt{txamg1110m000}.

Le descrizioni, numerazioni dei campi (\texttt{txamg.fdno}) ed i tipi di dominio correlati (\texttt{txamg.atyp}) vengono
riversati nella tabella di supporto \texttt{txamg111} e sono quindi disponibili e precompilati per le successive attività
di data-mapping manuale verso le tabelle target CloudSuite.

\section{Implementazione del Generatore di Codice}
\subsection{Template Engine e Creazione Dinamica}
La generazione dinamica del codice è gestita mediante apposite librerie DLL (\textit{Dynamic Link Library}) fornite
nel sistema Infor, quali \texttt{txamgdll0012} (Script Generator - Load dll) e \texttt{txamgdll0013} (Script Generator - Process dll).

Il codice sorgente Baan C viene generato riga per riga, definendo le firme delle funzioni
(es. \texttt{function extern long process.import.file...()}) e le variabili necessarie.
Per farlo in modo conforme all'infrastruttura di Infor LN, il generatore interroga direttamente il Data Dictionary del
sistema tramite le tabelle di sistema \texttt{ttadv} (in particolare \texttt{tttadv420} per le definizioni delle
tabelle, \texttt{tttadv422} per i campi e \texttt{tttadv400} per i domini).
Questa interrogazione è delegata al componente \texttt{txamgdll0010} (DD Integration Library), che funge
da strato di astrazione estraendo in automatico tipi di dato, lunghezze massime e array per assicurare che lo script
rispetti rigorosamente le definizioni fisiche.

Tali librerie servono ad analizzare la configurazione definita nelle tabelle di \textit{mapping} (ad esempio \texttt{txamg203})
per generare script eseguibili, riducendo drasticamente il tempo manuale richiesto per configurare le regole di trasformazione
standard e garantendo che il codice finale aderisca alla logica di business (\textit{business logic}).

\begin{listing}[H]
\begin{minted}[frame=lines, framesep=2mm, baselinestretch=1.2, fontsize=\footnotesize, linenos]{c}
function void GeneraMap(domain tcmcs.str8 tabella)
{
    Riga("function long map."&tabella&".record() {")
    select txamg203.*, txamg111.* from txamg203, txamg111
    where txamg203._index1 =
    {:txamg202.cbat, :txamg202.sequ, :txamg202.sseq}
    and txamg203.cmbd refers to txamg111
    selectdo
        on case txamg203.type
            case txamg.type.long:
                Riga("    "&tabella&trim$(txamg203.fdnm)
                          &" = lval(imp.str."&trim$(txamg203.fdnm)&")")
                break
            case txamg.type.string:
                Riga("    tt.align.according.domain(imp.str."
                         &trim$(txamg111.fdnm)&","
                         &tabella&trim$(txamg203.fdnm)&",...)")
                break
            |* Altri tipi (date, enum, etc.)
        endcase
    endselect
    Riga("    return(OK)")
    Riga("}")
}
\end{minted}
\caption{Logica di generazione della funzione di mapping (da \texttt{txamgdll0013}).}
\label{lst:mapping_gen}
\end{listing}

Nel seguente esempio viene mostrata la funzione \texttt{GeneraMap} che genera il codice per mappare i dati estratti dal
file alla struttura della tabella target:
\begin{itemize}
    \item \textbf{Riga 4-6:} Esegue un ciclo sui campi mappati (\texttt{txamg203}) relazionandoli con i campi sorgente (\texttt{txamg111}).
    \item \textbf{Riga 9-10:} Se il target è un numero intero (\textit{long}), genera una chiamata a \texttt{lval()} per convertire la stringa ASCII in valore numerico.
    \item \textbf{Riga 12-13:} Se il target è una stringa, usa \texttt{tt.align.according.domain()} per gestire correttamente lunghezze e allineamenti secondo il dominio Infor LN.
\end{itemize}

\subsection{Supporto alle Funzioni di Trasformazione Ad-hoc}
Oltre alle regole predefinite, l'interprete ammette la dichiarazione e l'aggancio di funzioni custom.
Durante il processo, eventuali routine di validazione o pulizia del dato possono essere inoltrate ad altre
interfacce di esecuzione (come \texttt{txamg3200} - Execute Batch Runs).

L'ausilio del componente \texttt{txamgdll0011} è importante in quanto si occupa del management dei metadati e dei log
contenenti eventuali errori (Error Logging).

\section{Interfacce Utente (Sessioni)}
L'interazione con l'utente avviene attraverso un gruppo di sessioni dedicate, denominate \textit{m-sessions} o
\textit{s-sessions} a seconda della loro finalità operativa.

Si evidenziano, per esempio:
\begin{itemize}
    \item \textbf{Setup iniziale:} la Sessione \texttt{txamg0100m000} permette la parametrizzazione globale dello strumento.
    \item \textbf{Definizione Formati:} le Sessioni relative a \texttt{txamg1110} gestiscono il popolamento dei formati dati elaborati.
    \item \textbf{Mapping Dati:} la definizione di conversione logica e associazioni campo-campo tra origine e
    destinazione avviene nelle interfacce di visualizzazione delle tabelle \textit{target} (\texttt{txamg2102m000}) e \textit{fields mapping} (\texttt{txamg2103m000}).
\end{itemize}

Nelle seguenti sezioni vengono riportati gli elenchi dei componenti software sviluppati, suddivisi per tipologia.

\subsection{Librerie (DLL) e Data Access Layer (DAL)}

\begin{table}[H]
    \centering
    \begin{tabular}{l l p{8cm}}
        \toprule
        \textbf{Tipo} & \textbf{Codice} & \textbf{Descrizione} \\
        \midrule
        Libreria & \texttt{txamgdll0010} & AlteaIN Migration Tool - DD Integration Library \\
        Libreria & \texttt{txamgdll0011} & AlteaIN Migration Tool - Metadata Management \& Error Logging \\
        Libreria & \texttt{txamgdll0012} & AlteaIN Migration Tool - Script Generator - Load DLL \\
        Libreria & \texttt{txamgdll0013} & AlteaIN Migration Tool - Script Generator - Process DLL \\
        \addlinespace
        DAL & \texttt{txamg000} & DAL for txamg000 \\
        DAL & \texttt{txamg111} & DAL for txamg111 \\
        DAL & \texttt{txamg203} & DAL for txamg203 \\
        DAL & \texttt{txamg1110} & Data Files \\
        DAL & \texttt{txamg3200} & Execute Batch Runs \\
        DAL & \texttt{txamg3500} & Runs by Batch \\
        DAL & \texttt{txamg3501} & Runs by Batch/Datafile \\
        DAL & \texttt{txamg3502} & Runs by Batch/Datafile/Table \\
        DAL & \texttt{txamg3503} & Runs Errors \\
        \bottomrule
    \end{tabular}
    \caption{Elenco delle Librerie (DLL) e DAL sviluppati.}
    \label{tab:dll_dal}
\end{table}

\subsection{Tabelle di Sistema e Sessioni Operative}

Elenco delle principali tabelle (le strutture dettagliate sono consultabili nell'Appendice \ref{appendix:tabelle}) e
delle sessioni utilizzate.

\begin{table}[H]
    \centering
    \begin{tabular}{l p{9cm}}
        \toprule
        \textbf{Tabella} & \textbf{Descrizione} \\
        \midrule
        \texttt{txamg000} & AlteaIN Migration Tool Parameters \\
        \texttt{txamg100} & Table groups definitions \\
        \texttt{txamg110} & Datafiles source definitions \\
        \texttt{txamg111} & DataFile Fields specifications \\
        \texttt{txamg200} & Batch definitions \\
        \texttt{txamg201} & Batch Sequences execution order \\
        \texttt{txamg202} & Target Tables mapping \\
        \texttt{txamg203} & MappedFields specific logic \\
        \texttt{txamg300} & Runs by Batch execution logs \\
        \texttt{txamg301} & Runs by Batch/Datafile detailed logs \\
        \texttt{txamg302} & Runs by Batch/Datafile/Table detailed logs \\
        \texttt{txamg303} & Runs Errors repository \\
        \bottomrule
    \end{tabular}
    \caption{Riepilogo delle tabelle introdotte dal Migration Tool.}
    \label{tab:tabelle}
\end{table}

\begin{table}[H]
    \centering
    \begin{tabular}{l p{9cm}}
        \toprule
        \textbf{Sessione} & \textbf{Descrizione} \\
        \midrule
        \texttt{txamg0100m000} & Parameters \\
        \texttt{txamg1100m000} & Datafile Groups \\
        \texttt{txamg1110m000} & Data Files \\
        \texttt{txamg1111m000} & DataFile Fields \\
        \texttt{txamg2100m000} & Batches \\
        \texttt{txamg2101m000} & Batch Sequences \\
        \texttt{txamg2102m000} & Batch Target Tables \\
        \texttt{txamg2103m000} & Table Mappings \\
        \texttt{txamg3200m000} & Execute Batch Runs \\
        \texttt{txamg3500m000} & Runs by Batch \\
        \texttt{txamg3503m000} & Runs Errors \\
        \bottomrule
    \end{tabular}
    \caption{Sessioni principali sviluppate per il Migration Tool.}
    \label{tab:sessioni}
\end{table}

\begin{table}[H]
    \centering
    \begin{tabular}{l p{9cm}}
        \toprule
        \textbf{Dominio} & \textbf{Descrizione} \\
        \midrule
        \texttt{txamg.atyp} & DataFile Field Type (Enumerated) \\
        \texttt{txamg.cbat} & Batch (String) \\
        \texttt{txamg.cdom} & Domain (String) \\
        \texttt{txamg.cmod} & Module (String) \\
        \texttt{txamg.cpac} & Package (String) \\
        \texttt{txamg.cprs} & Script/Library Code (String) \\
        \texttt{txamg.dfil} & Data File (String) \\
        \texttt{txamg.dfnm} & Data File Field Name (String) \\
        \texttt{txamg.fast} & ASCII File Format (Enumerated) \\
        \texttt{txamg.fdnm} & Field Name (String) \\
        \texttt{txamg.fdno} & Field Number (Integer) \\
        \texttt{txamg.flno} & Table (String) \\
        \texttt{txamg.grou} & Table Groups (String) \\
        \texttt{txamg.logn} & Login (String) \\
        \texttt{txamg.mapl} & Map Logic (Enumerated) \\
        \texttt{txamg.migt} & Migration Type (Enumerated) \\
        \texttt{txamg.refm} & Reference Mode (Enumerated) \\
        \texttt{txamg.type} & Domain Type (Enumerated) \\
        \texttt{txamg.utyp} & Update Type (Enumerated) \\
        \bottomrule
    \end{tabular}
    \caption{Domini personalizzati tx del Migration Tool.}
    \label{tab:domini}
\end{table}

\section{Integrazione Documentale}
\subsection{Esportazione delle Specifiche di Migrazione}
L'integrazione formale del modulo d'implementazione sfocia nella sua interconnessione con lo strumento aziendale adibito
alla documentazione del progetto.

Prima della configurazione ed eventuale esecuzione esplorativa del batch (le sessioni
\texttt{txamg2100} e \texttt{txamg3200}), l'\textit{utility} estrae il risultato dei mapping salvato nel \texttt{txamg203} (Tabella Mappings),
formattandolo utilizzando layout predefiniti.

L'esportazione documenta in PDF ogni singolo dominio trasferito, documentando sia
i campi copiati senza modifiche (\textit{one-to-one}) (in accordo con \texttt{txamg.mapl}) sia quei casi
particolari per cui è stato inserito del codice personalizzato, semplificando una possibile manutenzione logica dell'architettura in futuro.

\section{Ciclo di Vita del Dato}
Il processo di migrazione segue un flusso lineare, descritto nel seguente diagramma:

\begin{figure}[H]
    \centering
    \begin{tikzpicture}[node distance=2cm]
        \node (legacy) [startstop] {Legacy Data (ASCII)};
        \node (extract) [process, below of=legacy] {Estrazione};
        \node (parse) [process, below of=extract] {Parsing (Importer)};
        \node (map) [process, below of=parse] {Mapping \& Trasformazione};
        \node (gen) [process, below of=map] {Generazione Codice DLL};
        \node (exec) [process, below of=gen] {Esecuzione Batch};
        \node (db) [io, below of=exec] {Infor CloudSuite DB};
        \node (error) [startstop, right of=exec, node distance=5cm] {Log Errori (txamg303)};

        \draw [arrow] (legacy) -- (extract);
        \draw [arrow] (extract) -- (parse);
        \draw [arrow] (parse) -- (map);
        \draw [arrow] (map) -- (gen);
        \draw [arrow] (gen) -- (exec);
        \draw [arrow] (exec) -- (db);
        \draw [arrow] (exec) -- (error);
        \draw [arrow] (parse) -- (error);
    \end{tikzpicture}
    \caption{Diagramma di Flusso del Ciclo di Vita del Dato.}
    \label{fig:flowchart}
\end{figure}

Il flusso garantisce che ogni record sia tracciato.
In caso di errore durante il parsing (formato scorretto) o durante l'esecuzione (vincoli di database), l'engine
scrive un record dettagliato nella tabella \texttt{txamg303}, permettendo la ripartenza selettiva della migrazione.
